\section{Finite local boundary costs and the remaining growth question}
\label{sec:local-boundary-cost-scope}

\paragraph{Status and purpose.}
This scope clarification of MA-MO-036 has passed the complete
post-writing author reading and two complete independent readings.
It supplies no new arithmetic bound.  Meromorphicity removes strong local
convergence as a separate condition, but does not remove the
polynomial height bound.

Keep $0<a<1/2$, $s\ge0$, $r=s/2$, the structural finite Euler
factors, and the exact original source settings of MA-MO-036.
The rational sector retains both actual polar coordinates.
Use the same $\epsilon_0$ and uniform smooth partition $\chi_n$.
When the completed divisors lie in the closed strip, define the
actual local costs, allowing the value $+\infty$,
\[
 C_{\omega,n}
   =\sup_{0<\epsilon\le\epsilon_0}
       \left\|\frac{\chi_n(u)}
                       {A_\omega(a+\epsilon+iu)}\right\|_{H^{-r}} .
 \tag{MA-MO.1528}
\]
There are two right-boundary families for $\eta,\bar\eta$ and
one self-dual family for the rational principal case.  The
positive central strip excludes finite Euler zeros and the
rational pole exactly as in (MA-MO.1500).  No endpoint
$a=1/2$ assertion is made.

\paragraph{Finite-Laurent clarification of MA-MO-036.}
Under the closed-strip premise, every $C_{\omega,n}$ is finite
if and only if every completed boundary zero is fully retained:
its multiplicity $m$ satisfies $m<r+1/2$.
When these local costs are finite, each strong limit in
(MA-MO.1502) exists automatically.
Consequently, in MA-MO-036 its condition (L) can be replaced by
the single uniform estimate
$C_{\omega,n}\le C(2+|n|)^D$.
Finiteness for each $n$ separately is not this estimate.

\begin{proof}
Fix a compact frequency window.  The actual reciprocal is
meromorphic in a neighborhood of the corresponding compact
vertical rectangle.  It has only finitely many poles there.
Its right-edge poles are exactly the completed boundary zeros,
with their actual orders.  Enlarge the frequency window slightly,
remove the finitely many right-edge principal parts, and use a
sufficiently small neighborhood on the left so that no remaining
interior pole meets that neighborhood.  This choice is possible
for this fixed window by isolation of zeros; its size is not
claimed uniform in height.  On the support of a compact smooth
cutoff $\chi$, one therefore has
\[
 \frac{\chi(u)}{A_\omega(a+\epsilon+iu)}
   =\chi(u)\sum_{\gamma}\sum_{j=1}^{m_\gamma}
          \frac{c_{\gamma,j}}{(\epsilon+i(u-\gamma))^j}
                    +\chi(u)H_\epsilon(u).
 \tag{MA-MO.1529}
\]
Here $c_{\gamma,m_\gamma}\ne0$ are the actual Laurent
coefficients, and $H_\epsilon$ tends smoothly on this compact
window to $H_0$ as $\epsilon\downarrow0$.  Its localized
smooth seminorms are uniformly bounded for this window.
If necessary treat $\epsilon$ bounded away from zero separately
by compactness.  None of these constants is estimated in $n$.

For $j\ge1$ the unlocalized pole has the exact Laplace and
Sobolev formulas
\[
\begin{aligned}
 (\epsilon+i(u-\gamma))^{-j}
   &=\frac1{(j-1)!}\int_0^\infty
         t^{j-1}e^{-\epsilon t}e^{-i(u-\gamma)t}\,dt,\\
 \|(\epsilon+i(\,\cdot-\gamma))^{-j}\|_{H^{-r}}^2
   &=\frac{2\pi}{((j-1)!)^2}
       \int_0^\infty(1+t^2)^{-r}
                        t^{2j-2}e^{-2\epsilon t}\,dt .
\end{aligned}
\tag{MA-MO.1530}
\]
Thus if $j<r+1/2$, dominated convergence gives its strong
$H^{-r}$ right-boundary limit and a uniform norm bound.
Multiplication by $\chi$ is bounded on that Sobolev space.
If every boundary order is retained, the finite sum
(MA-MO.1529) and its smooth remainder therefore converge
strongly and have bounded norms.  Apply this to each $\chi_n$.
The finite number of poles in one window can have arbitrarily
close but distinct ordinates; no separation estimate is needed
for this assertion about a single window.

Conversely suppose every cost in (MA-MO.1528) is finite, and
let $a+i\gamma$ be a boundary zero of order $m$.
Since the partition sums to one, some $\chi_n(\gamma)\ne0$.
Choose $\psi\in C_c^\infty$, equal to one near $\gamma$,
supported where that $\chi_n$ does not vanish and avoiding
other boundary poles.  Multiplication by $\psi/\chi_n$
then gives a uniform $H^{-r}$ bound for $\psi/A_\omega$.
Its Laurent expansion is
$\psi\sum_{j=1}^m c_j(\epsilon+i(u-\gamma))^{-j}$
plus a uniformly smooth compact remainder.
Replacing this localized pole sum by the unlocalized one
changes it by a uniformly bounded $L^2$ function:
outside a fixed neighborhood of $\gamma$ the denominators
are bounded away from zero, and their tails are $O(|u|^{-1})$
or better.  Since $r\ge0$, this error is also bounded in
$H^{-r}$.

The unlocalized pole sum has Laplace inverse
\[
 \mathbf1_{t>0}e^{-\epsilon t}e^{i\gamma t}P(t),
 \qquad
 P(t)=\sum_{j=1}^m\frac{c_jt^{j-1}}{(j-1)!},
                       \qquad c_m\ne0 .
 \tag{MA-MO.1531}
\]
For all sufficiently large $t$, $|P(t)|\ge c t^{m-1}$.
Its squared Sobolev norm is therefore bounded below by
a positive constant times
$\int_{T_0}^\infty t^{2m-2-2r}e^{-2\epsilon t}\,dt$.
This diverges as $\epsilon\downarrow0$ when
$m\ge r+1/2$, logarithmically at equality.  That contradicts
the uniform localized bound.  Thus every right-boundary
order is retained.  Apply the same argument to the dual
sector; the functional equation gives the left boundary.
In the rational case self-duality gives the same conclusion.

We have proved that local finiteness forces retention,
and retention supplies strong convergence for every window.
A polynomial estimate on the costs therefore supplies both
parts of (MA-MO.1502), with their actual boundary values.
The converse follows immediately from that condition.
\end{proof}

\paragraph{Closed-strip containment can also be made a conclusion.}
Use exactly the universal bounded family (B) of MA-MO-036:
the two-sided fixed support for a nontrivial character, or
the positive-test actual negative columns into $\mathcal W$
for the rational case.  Its boundedness for any one nonzero
test already rules out a strict exterior zero.
For the nontrivial sector this is the genuine-core
obstruction (MA-MO.554)--(MA-MO.555); the dual relation
supplies the corresponding divisor statement.
For the rational negative column it is the actual
negative-source sequence (MA-MO.652), with both output
polar coordinates retained.  The zeta functional equation
supplies the reflected edge.  Thus one may state the exact
logical closure without a preliminary containment assumption:
\[
\begin{gathered}
 \text{the specified universal actual bounded family}\\
 \Longleftrightarrow
 \left\{
 \begin{array}{l}
 \text{all completed zeros lie in }|\Re(\rho-1/2)|\le a,\\
 C_{\omega,n}\le C(2+|n|)^D
                 \text{ for all the indicated right-edge sectors}
 \end{array}\right.\\
 \Longleftrightarrow
 \text{the full original matrix is trace class for every compact
 smooth test}.
\end{gathered}
\tag{MA-MO.1532}
\]
Indeed the first implication gives containment by the
exterior obstruction and then the costs by MA-MO-036.
The reverse implication and the last equivalence are
MA-MO-036 with the finite-Laurent clarification above.
The full retention and trace-norm continuity conclusions
are unchanged.  This is the same equivalent analytic gate,
with its hypotheses reorganized; no arithmetic decision
or additional movement is claimed.

\paragraph{What is still uncontrolled.}
Individual retention proves $C_{\omega,n}<\infty$ for each
fixed $n$.  The Laurent coefficients, the distances to
nearby interior zeros, and the analytic remainder constants
in (MA-MO.1529) may depend on that window.  Meromorphic
isolation gives no height estimate for them.
The already checked finite-disk argument bounds the
zero-free canceled factor polynomially; it leaves the
nearby pole product to be estimated, as in
(MA-MO.1414)--(MA-MO.1416).

For clarity, even the elementary local model with all poles
strictly inside the edge has, for any fixed integer
$k>p=r+1/2$,
\[
 \left\|(d+i\,\cdot\,)^{-k}\right\|_{H^{-r}}
                   \asymp_{k,r} d^{\,p-k}
                          \qquad(d\downarrow0).
\tag{MA-MO.1533}
\]
The upper bound is (MA-MO.1409); for the lower bound
restrict (MA-MO.1530) to $1/d\le t\le2/d$.
Multiplication by a fixed cutoff equal to one near zero
has the same divergent order, since the removed tails
are uniformly $L^2$.  This explains why local finiteness
cannot be substituted into the proof as a uniform bound.
It is a local analytic model, not a claimed configuration
of any actual Hecke divisor.  The conditional genuine-source
cluster obstructions MA-MO-033/034 remain explicitly
conditional on their arithmetic configurations.

For a fixed divisor let $\theta$ be the supremum of the
absolute centered real parts.  If $a>\theta$, the previously
proved fixed-gap theorem supplies the costs.  If $a<\theta$,
the exterior obstruction applies.  Within $0<a<1/2$, the
remaining no-margin case is therefore the extremal width
$a=\theta$, when that width exists in this interval.
This observation establishes no value of $\theta$ and
does not settle the polynomial cost there.

\paragraph{Primary-source check of the proposed arithmetic shortcut.}
The following sources were inspected for the missing estimate;
none is imported into the proof above.
Ng (2004), introduction and Theorem 3, treats Möbius
summatory behavior under RH together with an additional
upper bound for the negative derivative moment $J_{-1}$.
This separates a zero-location assumption from that
quantitative reciprocal information.

Bui, Florea, and Milinovich (2024), Sections 1.1--1.2 and
the complete proof of Lemma 2.1, obtain their stated
negative derivative-moment bounds on subfamilies with
specified small-gap control.  Their discussion of the
full family retains additional assumptions.  These results
do not furnish the present all-window Sobolev cost bound.
This comparison is not an assertion that their moment
problem is equivalent to (MA-MO.1528).

The more recent Bui, Florea, and Heycock preprint
(2026, version 1), Theorems 1.1--1.5 on printed pages 2--4,
assumes RH and addresses specified positive shifts,
averages, or typical short intervals.  Those statement
ranges do not supply a uniform $\epsilon\downarrow0$
bound in every fixed window here.  Only these statements
and their hypotheses were checked; its complete proof
was not re-audited or used as an input.

Accordingly this investigation has not established that
classical estimates or the actual Hecke arithmetic force
polynomial growth of (MA-MO.1528) from individual retention.
It also has not proved that such growth fails for an actual
Hecke function.  A supplies the already separated rational
and Dirichlet fixed-gap branch; M supplies the specified
Hecke branch.  The independent retention-only gate remains
open, and FT002 has not moved.

\paragraph{References and analytic placement.}
Ng, N. (2004). The distribution of the summatory function
of the Möbius function. \emph{Proceedings of the London
Mathematical Society, 89}(2), 361--389.
\url{https://doi.org/10.1112/S0024611504014741}.
Author manuscript:
\url{https://www.cs.uleth.ca/~nathanng/RESEARCH/mobius2b.pdf}.

Bui, H. M., Florea, A., \& Milinovich, M. B. (2024).
Negative discrete moments of the derivative of the
Riemann zeta-function. \emph{Bulletin of the London
Mathematical Society, 56}(8), 2680--2703.
\url{https://doi.org/10.1112/blms.13092}.

Bui, H. M., Florea, A., \& Heycock, A. (2026).
\emph{Small values of the Riemann zeta-function near the
critical line} (Version 1). arXiv.
\url{https://arxiv.org/abs/2610.27591v1}.

The finite-length divisor modules and their exterior-power
determinants retain the inherited cote 1 placement of
MA-MO-036, BR-TR-01/02, scans 079--080 and 111--112.
The present clarification concerns analytic norms and
their quantifiers; it adds no positivity or arithmetic
consequence to those finite-length $K_0$ classes.
